\section{总结与延伸}

\begin{enumerate}
    \item 行为克隆提供最快的入门，但是需要警惕均值行为与 covariate shift。
    \item 生成模型使策略具备多模态表达力，需要在采样阶段做 KL/temperature 调整。
    \item DAgger/HG-DAgger/混合 RL 策略是解决 compounding errors 的有效路径。
    \item 高质量演示数据的 audit、metadata、multi-role 协同是落地的基础。
    \item 部署环节要以 evidence board 为中心，把 benchmark、safety、ops 三条线绑定到监控迭代。
    \item 记录每次 mitigation 与 rollback 事件，让 governance audit 有据可依。
\end{enumerate}

\subsection{落地路线图}

\begin{figure}[H]
\centering
\includegraphics[width=0.8\textwidth]{slide-02-14.jpg}
\caption{Slide 02-14 以 timeline 形式展示了从数据到 deployment 的 deliverable。}
\end{figure}

落地路线图将模仿学习 pipeline 拆成三个阶段：1) 数据准备（demo catalog + metadata），2) 策略训练（diffusion/MoG + flow calibration），3) 灾难恢复（drift playbook + governance review）。每个阶段都要定义 guardrail、evaluation metric 与 artifacts。

\begin{table}[H]
\centering
\begin{tabular}{p{0.25\textwidth} p{0.55\textwidth}}
\toprule
阶段 & Key artifact \\
\midrule
数据准备 & Metadata-rich demo catalog + coverage report \\
策略训练 & KL drift sweep + sampling temperature record \\
部署恢复 & Drift playbook + rollback checklist \\
\bottomrule
\end{tabular}
\caption{模仿学习落地的阶段化路线图}
\end{table}

\begin{warningbox}{别把 roadmap 当计划书}
Roadmap 应该不断更新，从 evidence dashboard 获取实际 delay 和 drift，再把 lesson learned 反向写回 timeline，才能持续前进。
\end{warningbox}

\subsection{总结表}

\begin{table}[H]
\centering
\begin{tabular}{p{0.2\textwidth} p{0.35\textwidth} p{0.35\textwidth}}
\toprule
主题 & 核心 takeaway & 工程行动 \\
\midrule
行为克隆 & 快速上线但需注意 covariance shift & 训练后跑 KL sweep、部署前做 sanity check \\
策略表达 & 生成分布需要被采样探索 & 使用 diffusion/MoG 并记录 sampling temperature \\
在线修正 & DAgger 弥补未见状态 & 建立 Observe-Assess-Act-Document runbook \\
数据 QA & Multi-role audit & 用 metadata + clustering 识别 coverage gaps \\
部署治理 & Evidence dashboard & 绑定 benchmark/safety/ops 3 条线 \\
\bottomrule
\end{tabular}
\caption{本讲的行动化提炼}
\end{table}

\subsection{拓展阅读}
\begin{itemize}
    \item Chi et al., \textit{Diffusion Policy}：将 diffusion model 引入机器人策略。
    \item Ross et al., \textit{A Reduction of Imitation Learning and Structured Prediction to No-Regret Online Learning}：DAgger 原文。
    \item Brohan et al., \textit{RT-2}：展示多模态模仿学习的效果。
    \item Rajeswaran et al., \textit{Behavioral Cloning with Expert Interventions}：将 human interventions 系统化。
\end{itemize}

\subsection{本章小结}

本讲从行为克隆到生成策略、从 DAgger/online pipeline 到 data QA 与部署治理，形成一个 end-to-end 的模仿学习实践框架。只有把各个阶段都打通，才能让模仿学习在生产环境中真正可靠地执行。

\end{document}
